\documentclass[border=10pt]{standalone}
\usepackage[no-math]{fontspec}
\usepackage[brazil]{babel}
\usepackage{amsmath,amssymb}
\IfFileExists{C:/Windows/Fonts/Lato_Regular.ttf}{%
  \setmainfont{Lato}[
    Path = C:/Windows/Fonts/,
    UprightFont = Lato_Regular.ttf,
    BoldFont = Lato_Bold.ttf,
    ItalicFont = Lato_Italic.ttf,
    BoldItalicFont = Lato_BoldItalic.ttf
  ]
  \setsansfont{Lato}[
    Path = C:/Windows/Fonts/,
    UprightFont = Lato_Regular.ttf,
    BoldFont = Lato_Bold.ttf,
    ItalicFont = Lato_Italic.ttf,
    BoldItalicFont = Lato_BoldItalic.ttf
  ]
}{%
  \RequirePackage[default]{lato}
}
\renewcommand{\familydefault}{\sfdefault}

\usepackage{xcolor}
\usepackage{tikz}
\usepackage{pgfplots}
\pgfplotsset{compat=1.18}
\usetikzlibrary{arrows.meta,positioning,shapes.geometric,fit,backgrounds,calc,decorations.pathreplacing,patterns}

% Paleta institucional UEFS / Sociologia Rural
\definecolor{srNavy}{HTML}{1F4268}
\definecolor{srNavyDark}{HTML}{152E4A}
\definecolor{srGold}{HTML}{B38A3D}
\definecolor{srGoldDark}{HTML}{8C6927}
\definecolor{srTeal}{HTML}{0E7C7B}
\definecolor{srForest}{HTML}{2D6A4F}
\definecolor{srAlert}{HTML}{9E2A2B}
\definecolor{srInk}{HTML}{26313D}
\definecolor{srMuted}{HTML}{5B6573}

\definecolor{bgNavy}{HTML}{EDF2F7}
\definecolor{bgGold}{HTML}{F8F3E9}
\definecolor{bgTeal}{HTML}{EBF6F5}
\definecolor{bgForest}{HTML}{EEF5F0}
\definecolor{bgAlert}{HTML}{FDF2F2}

\tikzset{
  >=Stealth,
  srBox/.style={
    draw=#1, fill=white, rounded corners=5pt, line width=1.15pt,
    inner sep=7pt, align=center, font=\footnotesize\sffamily, text=srInk
  },
  srTint/.style={
    draw=#1, fill=#1!8!white, rounded corners=5pt, line width=1.15pt,
    inner sep=7pt, align=center, font=\footnotesize\sffamily, text=srInk
  },
  srBadge/.style={
    fill=#1, text=white, rounded corners=9pt, inner xsep=8pt, inner ysep=3pt,
    font=\scriptsize\bfseries\sffamily
  },
  srTitle/.style={font=\normalsize\bfseries\sffamily, text=srNavy, align=center},
  srSub/.style={font=\footnotesize\sffamily, text=srMuted, align=center},
  srArrow/.style={->, >=Stealth, line width=1.35pt},
  srDArrow/.style={<->, >=Stealth, line width=1.35pt}
}


\begin{document}
\begin{tikzpicture}
  % Título
  \node[srTitle, font=\large\bfseries\sffamily] (title) at (0, 3.05)
    {Movimentos Sociais do Campo e Repertórios de Ação no Brasil};
  \node[srSub] (sub) at (0, 2.52)
    {Trajetórias históricas, sujeitos políticos e formas de organização no meio rural};

  \def\pW{3.55cm}
  \def\pH{4.85cm}
  \def\hdrW{3.2cm}
  \def\bodyW{3.1cm}

  % Pilar 1: Ligas Camponesas
  \node[draw=srAlert, fill=bgAlert, rounded corners=6pt, line width=1.2pt,
        minimum width=\pW, minimum height=\pH, anchor=north] (p1) at (-5.85, 1.95) {};
  \node[fill=srAlert, text=white, rounded corners=4pt, text width=\hdrW,
        inner ysep=5pt, font=\footnotesize\bfseries\sffamily, align=center, anchor=north]
        (h1) at ($(p1.north)+(0,-0.15)$) {Ligas Camponesas\\(1955--64)};
  \node[text=srInk, font=\footnotesize\sffamily, text width=\bodyW, align=left, anchor=north]
        at ($(h1.south)+(0,-0.22)$) {%
    $\bullet$\;\;Liderança de Francisco Julião (Nordeste)\par\vspace{6pt}
    $\bullet$\;\;Luta pela Reforma Agrária ``na lei ou na marra''\par\vspace{6pt}
    $\bullet$\;\;Combate ao cambão e ao latifúndio tradicional
  };

  % Pilar 2: Sindicalismo & CONTAG
  \node[draw=srNavy, fill=bgNavy, rounded corners=6pt, line width=1.2pt,
        minimum width=\pW, minimum height=\pH, anchor=north] (p2) at (-1.95, 1.95) {};
  \node[fill=srNavy, text=white, rounded corners=4pt, text width=\hdrW,
        inner ysep=5pt, font=\footnotesize\bfseries\sffamily, align=center, anchor=north]
        (h2) at ($(p2.north)+(0,-0.15)$) {Sindicalismo \&\\CONTAG};
  \node[text=srInk, font=\footnotesize\sffamily, text width=\bodyW, align=left, anchor=north]
        at ($(h2.south)+(0,-0.22)$) {%
    $\bullet$\;\;Direitos trabalhistas e Estatuto do Trabalhador Rural\par\vspace{6pt}
    $\bullet$\;\;Previdência social rural (Funrural / CF 1988)\par\vspace{6pt}
    $\bullet$\;\;Negociação sindical e políticas de crédito
  };

  % Pilar 3: MST, CPT & Reforma Agrária
  \node[draw=srForest, fill=bgForest, rounded corners=6pt, line width=1.2pt,
        minimum width=\pW, minimum height=\pH, anchor=north] (p3) at (1.95, 1.95) {};
  \node[fill=srForest, text=white, rounded corners=4pt, text width=\hdrW,
        inner ysep=5pt, font=\footnotesize\bfseries\sffamily, align=center, anchor=north]
        (h3) at ($(p3.north)+(0,-0.15)$) {MST, CPT \&\\Reforma Agrária};
  \node[text=srInk, font=\footnotesize\sffamily, text width=\bodyW, align=left, anchor=north]
        at ($(h3.south)+(0,-0.22)$) {%
    $\bullet$\;\;Ocupações de terras improdutivas e acampamentos\par\vspace{6pt}
    $\bullet$\;\;Assentamentos, cooperativas e agroindústrias\par\vspace{6pt}
    $\bullet$\;\;Educação do campo e agroecologia
  };

  % Pilar 4: MAB, MMC & Via Campesina
  \node[draw=srTeal, fill=bgTeal, rounded corners=6pt, line width=1.2pt,
        minimum width=\pW, minimum height=\pH, anchor=north] (p4) at (5.85, 1.95) {};
  \node[fill=srTeal, text=white, rounded corners=4pt, text width=\hdrW,
        inner ysep=5pt, font=\footnotesize\bfseries\sffamily, align=center, anchor=north]
        (h4) at ($(p4.north)+(0,-0.15)$) {MAB, MMC \&\\Via Campesina};
  \node[text=srInk, font=\footnotesize\sffamily, text width=\bodyW, align=left, anchor=north]
        at ($(h4.south)+(0,-0.22)$) {%
    $\bullet$\;\;Atingidos por barragens (MAB) e justiça hídrica\par\vspace{6pt}
    $\bullet$\;\;Mulheres camponesas (MMC) e protagonismo feminino\par\vspace{6pt}
    $\bullet$\;\;Soberania alimentar e sementes crioulas
  };
\end{tikzpicture}
\end{document}
